\documentclass[10pt]{article}
\usepackage[a4paper,margin=0.7in]{geometry}
\usepackage{amsmath,amssymb,mathtools,booktabs,longtable,array}
\usepackage[dvipsnames]{xcolor}
\usepackage[colorlinks=true,linkcolor=blue,citecolor=blue,urlcolor=blue]{hyperref}
\allowdisplaybreaks
\numberwithin{equation}{section}
\newcommand{\JJ}{{\rm JIMWLK}\otimes{\rm LO}}
\newcommand{\BB}{{\rm BFKL}\otimes{\rm LO}}
\title{Our final result (mid rapidity) and AABKL, arXiv:1411.2869 (forward rapidity)}
\date{}
\begin{document}
\maketitle
\noindent Our result: \texttt{Remaining\_RC\_DGLAP\_finite.pdf}, Eq.~(8.1), with the evolution terms of \texttt{Checks\_JIMWLK\_BFKL.pdf}. AABKL: T.~Altinoluk, N.~Armesto, G.~Beuf, A.~Kovner, M.~Lublinsky, ``Single-inclusive particle production in proton-nucleus collisions at next-to-leading order in the hybrid formalism'', their final result Eq.~(4.22), the BK derivation Eqs.~(5.31)--(5.39), and App.~C. $a\equiv\frac{k^+}{\vee}$.

\section{Answer}
\begin{itemize}
\item \textbf{Term by term the two results do not coincide, and they should not.} The leading orders are different objects, so the NLO corrections are corrections to different things:
\begin{itemize}
\item AABKL (forward): a collinear parton of the proton, with a PDF $x_pf(x_p)$, scatters eikonally; LO is $O(\alpha_s^0)$; NLO is one splitting of that parton.
\item Ours (mid rapidity): a gluon radiated by the colour charges $\rho$ of a dilute projectile; LO is $O(\alpha_s)$ with $\langle\rho\rho\rangle$; NLO is one more gluon in $\Lambda<p^+<\vee$.
\end{itemize}
The finite constants are also scheme dependent: AABKL cut the soft gluon by an Ioffe time, we by sharp $p^+$ cutoffs.
\item \textbf{Everything that must agree does agree} (Sec.~\ref{sec:same}): the target evolution, the rapidity bookkeeping, the final-state DGLAP term (same $P_{gg}$, same plus prescription, same kinematic limit, same scheme constant), the hadron convolution, and the $\beta_0/2$ part of $P_{gg}$.
\item \textbf{The forward limit of our result has AABKL's structure} (Sec.~\ref{sec:forward}): at fixed $a=\frac{k^+}{\vee}=O(1)$ the BFKL logarithm is $O(1)$, and only the target evolves, as in their Eq.~(5.32).
\item \textbf{What is new at mid rapidity} (Sec.~\ref{sec:new}): the BFKL evolution of the projectile over $\log\frac{\vee}{k^+}=\frac12\log\frac{\vee}{\Lambda}$, which is as large as the target interval; the running of the LO coupling; and terms that exist only because the projectile is a set of colour charges.
\end{itemize}

\section{The two final results side by side}
\paragraph{Ours.}
\begin{align}
\frac{d^3N_h}{d^3k}={}&\int_{k^+/\vee}^1\frac{d\zeta}{\zeta^3}\,D_{h/g}(\zeta,\mu_F)\,\frac{d^3N_g}{d^3k}\Big(\frac{k^+}{\zeta},\frac{\mathbf k}{\zeta}\Big),\notag\\
\frac{d^3N_g}{d^3k}={}&{\rm LO}\big|_{\alpha_s(\mu_R)}+\frac{\alpha_s\beta_0}{4\pi}\log\frac{\mu_R^2}{\mu_F^2}\,{\rm LO}+\frac{\alpha_s}{2\pi}N_c\Big[\frac{67}{9}-\frac{\pi^2}{3}\Big]{\rm LO}
+\mathcal N\Big[\log\frac{k^+}{\Lambda}\,\JJ+\log\frac{\vee}{k^+}\,\BB\Big]\notag\\
&-\frac{\alpha_s}{2\pi}\int_{a}^{1}\frac{d\xi}{\xi^2}P^{(r)}_{gg,+}(\xi)\,{\rm LO}_\xi\big[\ell_\Delta(\mu_F)\big]-\log\frac{\vee}{k^+}H_4
+\frac{\alpha_sN_c}{4\pi}\log(1-a)\,{\rm LO}\big[\hat E_X+\hat E_{X'}\big]+\frac{d^3N_{V_\vee}}{d^3k}+\text{finite}.
\label{eq:ours}
\end{align}
\paragraph{AABKL, Eq.~(4.22)} (gluon channel; quark channels analogous), schematically:
\begin{align}
\frac{d\sigma}{d^2p_\perp d\eta}={}&\int_{x_F}^1\frac{d\zeta}{\zeta^2}D(\zeta)\Big\{x_pf^g_{\mu^2}(x_p)\,s^A_{Y_T}[y,\bar y]\ (\text{LO, }O(\alpha_s^0))\notag\\
&+\frac{g^2}{(2\pi)^3}\int_0^{1-x_p}d\xi\,\frac{x_p}{1-\xi}f^g\Big(\frac{x_p}{1-\xi}\Big)\Big[\frac{1-\xi}{[\xi]_+}+\frac{\xi}{1-\xi}+\xi(1-\xi)\Big]2C_A\notag\\
&\qquad\times\Big\{\big[A^i(y-z)A^i(\bar y-z)-(A^iA^i)_{\mu^2}\big]\Big(s_A[y,\bar y]\ \text{(PDF)}+(1-\xi)^2s_A[(1-\xi)y,(1-\xi)\bar y]\ \text{(FF)}\Big)+\dots\Big\}\notag\\
&+\text{real and virtual terms with }s\,s\,s,\ \frac{1}{N_c^2}X,\ s_A\ \text{(colour structures of the gluon channel)}\notag\\
&+\frac{g^2}{(2\pi)^3}x_pf^g(x_p)\int_0^1\frac{d\xi}{\xi}\int[A^i_\xi(y-z)-A^i_\xi(\bar y-z)]^2\Big[\frac{N_c^3}{N_c^2-1}(sss+sss)-\frac{1}{N_c^2}X-2C_As_A\Big]\Big\},
\label{eq:aabkl}
\end{align}
with the target evolved over $Y_T$, the projectile only over $Y_P=\ln\frac{1}{x_p}+O(1)$ (their Eq.~(2.11)), and the PDFs and FFs in their scheme, $f_{\mu^2}=f^{\overline{MS}}_{1.26\mu^2}$. The last line is the soft-gluon term with the Ioffe-time cutoff in $A_\xi$; they stress it is not an extra evolution, since its ``evolution interval'' depends on the momentum of the emitted gluon.

\section{What agrees}\label{sec:same}
{\small\renewcommand{\arraystretch}{1.3}
\begin{longtable}{@{}>{\raggedright\arraybackslash}p{0.27\textwidth}>{\raggedright\arraybackslash}p{0.30\textwidth}>{\raggedright\arraybackslash}p{0.36\textwidth}@{}}\toprule
AABKL & ours & how it was checked\\\midrule
target evolution: BK for the dipole, from $s_0\frac{\partial}{\partial s_0}\sigma=0$, Eqs.~(5.32)--(5.39) & $\log\frac{k^+}{\Lambda}\,\JJ$, from the independence of $\Lambda$ & our two-$\rho$ terms $\propto\log\frac{k^+}{\Lambda}$ are JIMWLK$\otimes$LO exactly (154 colour structures); BK is its large-$N_c$ dipole form. Same derivation: cutoff independence\\
projectile interval $Y_P=\ln\frac1{x_p}$, Eq.~(2.11) & $\log\frac{\vee}{k^+}$ & $x_p\leftrightarrow\frac{k^+}{\vee}$ (largest projectile momentum $\leftrightarrow P^+$)\\
FF (after-target) term: $P_{gg}$ with the plus prescription at the soft end, $\int_0^{1-x_p}d\xi$, $\log\big((y-\bar y)^2\mu^2\big)$, App.~C & DGLAP term $-\frac{\alpha_s}{2\pi}\int_a^1\frac{d\xi}{\xi^2}P^{(r)}_{gg,+}{\rm LO}_\xi[\ell_\Delta(\mu_F)]$ from Group II Row IV & with $\xi_{\rm AABKL}=1-\xi$: $P_{\rm ours}(\xi)-2N_cP_{\rm AABKL}(1-\xi)=0$; same limit ($1-x_p\leftrightarrow1-a$); $\ell_\Delta=\log\frac{(y'-y)^2\mu^2}{c_0^2}$ (\texttt{dglap\_rc/checks.json})\\
scheme: $f_{\mu^2}=f^{\overline{MS}}_{1.26\mu^2}$, $R=2e^{\Psi(1)}$, Eq.~(7.22) & $\overline{MS}$ with $c_0=2e^{-\gamma_E}$ & $R=c_0=1.12292$, $R^2=c_0^2=1.26095$ exactly (\texttt{checks.json})\\
$\frac{\beta_0}{2}\delta(1-\xi)$ of $P_{gg}$: from their virtual term, $\int_0^1d\xi\big[\frac{1-\xi}{[\xi]_+}+\frac{\xi}{[1-\xi]_+}+\xi(1-\xi)\big]=-\frac{11}{6}$, $2C_A\cdot\frac{11}{6}=\beta_0$ & from the coupling of LO (A3 and the shockwave through the loop cancel, \texttt{Steps\_line1\_and\_mu.pdf}) & both complete $P^{(r)}_{gg,+}$ to the full $P_{gg}$, which is absorbed in $D_{h/g}(\mu_F)$\\
hadron convolution $\int_{x_F}^1\frac{d\zeta}{\zeta^2}D(\zeta)$ & $\int_{k^+/\vee}^1\frac{d\zeta}{\zeta^3}D_{h/g}(\zeta)$ & same form\\
colour structures of the one-gluon real/virtual terms ($sss$, $X/N_c^2$, $s_A$) & ${\rm Tr}\{T^eAT^eB\}$ traces & same colour algebra; the kernels differ (emission by charges)\\\bottomrule
\end{longtable}}

\section{The forward limit of our result}\label{sec:forward}
AABKL (their Eq.~(5.31)): ``production at fixed $x_p$ but at higher energy''. For us that is $\vee\to\vee e^\Delta$ at fixed $a=\frac{k^+}{\vee}$, i.e.\ $k^+\to k^+e^\Delta$. Then $\log\frac{\vee}{k^+}=\log\frac1a$ is fixed and $\log\frac{k^+}{\Lambda}$ grows by $\Delta$. From Eq.~\eqref{eq:ours},
\begin{equation}
\frac{\partial}{\partial\log\vee}\frac{d^3N_g}{d^3k}\bigg|_{a}=\mathcal N\,\JJ+O(\alpha_s^2),
\end{equation}
only the target evolves. This is their Eq.~(5.32): the energy dependence at fixed $x_p$ is the evolution of the target ($s_{Y_T}$), BK in their large-$N_c$ dipole form, Eq.~(5.39). In this limit our $\log\frac{\vee}{k^+}\,\BB=\log\frac1a\,\BB$ is an $O(\alpha_s)$ term without a large logarithm, like their projectile evolution over $Y_P=\ln\frac1{x_p}=O(1)$, which they absorb into the PDF at the scale of the projectile (their Sec.~III.A). So at leading logarithm the forward limit of our result is AABKL's structure.

\section{What is different, and why}\label{sec:new}
\begin{itemize}
\item \textbf{BFKL of the projectile.} At mid rapidity $k^+=\sqrt{\vee\Lambda}$, so $\log\frac{\vee}{k^+}=\log\frac{k^+}{\Lambda}=\frac12\log\frac{\vee}{\Lambda}$: the projectile interval is as large as the target one and must be resummed. It is $\log\frac{\vee}{k^+}\,\BB$, your App.~E formula. AABKL say it themselves: ``For small values of $x$ one indeed has to take into account the $P^+$ dependence'' of the projectile distribution (Sec.~III.A, item 3). Forward it is absent because $\ln\frac1{x_p}=O(1)$.
\item \textbf{Running coupling.} Our LO is $\propto g^2$, so its coupling is renormalized: $\frac{\alpha_s\beta_0}{4\pi}\log\frac{\mu_R^2}{\mu_F^2}{\rm LO}$ and the constant $\frac{\alpha_s}{2\pi}N_c(\frac{67}{9}-\frac{\pi^2}{3}){\rm LO}$. AABKL's LO is $O(\alpha_s^0)$; they have no such term, and $\beta_0$ appears only in $P_{gg}$.
\item \textbf{No PDF DGLAP.} Our projectile is $\rho$, not a collinear parton: the parent gluon is emitted at distance $|x-w|$ from a charge, so the before-target emission (Group II Row I) is finite. AABKL have a PDF term $\propto P\otimes f\log\mu^2$.
\item \textbf{Terms that exist only for a projectile of colour charges:} the linear-polarization term ($\frac{k^ik^j}{k^2}$, line 3 of Group II Row IV), emission of the LO gluon after the target by the rotated charge, and the three- and four-$\rho$ terms.
\item \textbf{The limit $\xi\ge\frac{k^+}{\vee}$ in our virtual term} (Group I Row II) gives the pole Q1 $\propto\log(1-a)$ and $V_\vee=-\frac{\alpha_s}{2\pi}N_c\log(1-a)\log(\frac1a-1){\rm LO}$. AABKL's virtual term has $\int_0^1d\xi$ and neither term. Both vanish at mid rapidity in the high-energy limit ($a\to0$). With AABKL's choice $\int_0^1$ they would be absent from ours too.
\item \textbf{Quark channels} ($C_F$ terms, their Eq.~(4.21)): only in AABKL; we have $n_f=0$.
\item \textbf{Finite constants are scheme dependent:} the Ioffe-time cutoff vs.\ sharp cutoffs moves $O(\alpha_s)$ terms between ``evolution'' and ``finite'' (their Eqs.~(3.11)--(3.13) and their last term in Eq.~(4.22)). A numerical agreement of the finite parts is therefore not expected even where the structures are the same.
\end{itemize}

\section{Verdict}
Our result is consistent with AABKL where it has to be: the same target evolution and rapidity bookkeeping (forward limit = their structure), the same final-state DGLAP term in the same scheme, and the same $P_{gg}$ including its $\beta_0/2$ part. It differs where it must: at mid rapidity the projectile side has its own large logarithm (BFKL), our LO carries a coupling that runs, and a projectile of colour charges produces terms a collinear parton does not. A term-by-term match is not possible because the two leading orders, and the two soft-gluon regularizations, are different.
\end{document}
